% Manual · Laboratorio 02 — Ubuntu 24.04: estación receptora ADS-B/Mode S en 1090 MHz.

\PortadillaLaboratorio{Laboratorio 02 · Ubuntu 24.04 LTS · 225 minutos}{Del pulso de 1090 MHz a la aeronave en el mapa: recepción y decodificación ADS-B/Mode S}{Su misión: montar una estación pasiva que reciba los mensajes que emiten los aviones, demostrar que esos mensajes son válidos y llevarlos a una tabla y a un mapa. Terminará contrastando su estación con una grabación de contenido conocido.}

\section{Ficha del laboratorio}
\label{sec:u-plan}

\begin{longtable}{@{}>{\raggedright\arraybackslash}p{0.20\linewidth}>{\raggedright\arraybackslash}p{0.73\linewidth}@{}}
\toprule
\endhead
\textbf{Al terminar podrá} & preparar Ubuntu y verificar el HackRF y sus permisos; configurar una cadena de recepción en 1090 MHz; obtener y reconocer tramas Mode S y ADS-B válidas; interpretar una tabla de aeronaves y un mapa; explicar cómo se calcula una posición y por qué puede fallar; medir la calidad de la recepción con un diseño experimental honesto; y separar la información recibida, la calculada y la consultada en bases de datos. \\
\textbf{Antes de venir} & Responda el cuestionario diagnóstico (Anexo~B, p.~\pageref{B-sec:b-diagu}) y lea el \anexoA{sec:a-sistemas} y el \anexoA{sec:a-trama}. \\
\textbf{Equipo} & HackRF One, antena de 1090 MHz, filtro de 1090 MHz, adaptadores SMA, adaptador N -- SMA, SIGLENT SSA3032X Plus, PC con Ubuntu 24.04 y la grabación de referencia en su carpeta personal (sección~\ref{sec:banco}). \\
\textbf{Registro} & Hojas U1 a U11 del Anexo B. \\
\bottomrule
\end{longtable}

\begin{longtable}{@{}>{\raggedright\arraybackslash}p{0.62\linewidth}>{\raggedleft\arraybackslash}p{0.12\linewidth}>{\raggedleft\arraybackslash}p{0.16\linewidth}@{}}
\toprule
\textbf{Etapa} & \textbf{Minutos} & \textbf{Acumulado} \\
\midrule
\endhead
Orientación y etapa 1 · Montar la estación de 1090 MHz & 20 & 20 \\
2 · Comprobar Ubuntu & 20 & 40 \\
3 · Receptor y demodulador & 30 & 70 \\
4 · De la energía a la trama & 20 & 90 \\
5 · La tabla de aeronaves & 20 & 110 \\
6 · El mapa & 25 & 135 \\
7 · Medir la recepción & 25 & 160 \\
8 · El analizador en 1090 MHz & 20 & 180 \\
9 · Contraste con una grabación de referencia & 25 & 205 \\
10 · Cierre y entrega & 20 & 225 \\
\bottomrule
\end{longtable}

\begin{figure}[H]
\centering
\resizebox{\linewidth}{!}{% Laboratorio 02: arquitectura física y lógica de la estación receptora pasiva.
\begin{tikzpicture}[node distance=5mm and 7mm]
  \node[bloqueext,minimum width=27mm] (xpdr) {Transpondedor de la\\aeronave (Mode S,\\ADS-B Out 1090ES)};
  \node[bloque,right=22mm of xpdr,minimum width=18mm] (ant) {Antena\\receptora};
  \node[bloque,right=of ant,minimum width=21mm] (hrf) {\textbf{HackRF One}\\1090,000 MHz\\(solo RX)};
  \node[bloque,right=of hrf,minimum width=20mm] (pc) {Ubuntu 24.04\\USB + udev};
  \node[bloquealt,right=of pc,minimum width=23mm] (in) {SDRangel\\\textbf{HackRF Input}\\I/Q a 4 MSa/s};
  \draw[flecha,decorate,decoration={snake,amplitude=1.1mm,segment length=3mm,post length=2.5mm}] (xpdr) -- (ant);
  \node[rotulo] at ($(xpdr.east)!0.5!(ant.west)+(0,6mm)$) {1090 MHz\\ráfagas PPM};
  \draw[flecha] (ant) -- (hrf);
  \draw[flecha] (hrf) -- (pc);
  \draw[flecha] (pc) -- (in);

  \node[bloquealt,below=11mm of in,minimum width=27mm] (dem) {\textbf{ADS-B Demodulator}\\preámbulo, bits, paridad};
  \node[bloquealt,left=of dem,minimum width=23mm] (msg) {Tramas válidas\\(DF, ICAO, TC)};
  \node[bloquealt,left=of msg,minimum width=22mm] (tab) {Tabla de\\aeronaves};
  \node[bloquealt,left=of tab,minimum width=24mm] (map) {Feature \textbf{Map}\\(posición CPR)};
  \draw[flecha] (in) -- (dem); \draw[flecha] (dem) -- (msg); \draw[flecha] (msg) -- (tab); \draw[flecha] (tab) -- (map);

  \node[bloqueext,below=10mm of map,minimum width=30mm] (pos) {\menu{Preferences} $\rightarrow$ \textbf{My Position}\\(referencia CPR y distancia)};
  \node[bloqueext,below=10mm of msg,minimum width=30mm,xshift=-6mm] (db) {Bases externas descargadas\\(enriquecen; \emph{no} son RF)};
  \node[bloqueinst,below=10mm of dem,minimum width=25mm,xshift=8mm] (file) {\textbf{FileInput} (.sdriq)\\respaldo grabado};
  \draw[flechadash] (pos) -- (map);
  \draw[flechadash] (db.north) -- (tab.south east);
  \draw[flechadash] (file.north) -- ([xshift=8mm]dem.south);
\end{tikzpicture}
}
\caption{Arquitectura de la estación. Las flechas continuas siguen la señal; las discontinuas son entradas que \emph{no} vienen de la aeronave: la posición configurada, las bases de datos y, en la etapa 9, el archivo grabado.}
\label{fig:cadenaadsb}
\end{figure}

\begin{alerta}[Una estación que solo escucha]
Su estación no interroga a nadie. Recibe las difusiones ADS-B y las respuestas de los transpondedores a radares secundarios y a equipos anticolisión. Está prohibido emitir en 1030 y 1090 MHz (sección~\ref{sec:reglas}).
\end{alerta}

% ---------------------------------------------------------------------------
\section{Etapa 1 · Montar la estación de 1090 MHz}
\label{sec:u-antena}

\Puente{A 1090 MHz las señales viajan casi en línea recta y las paredes las atenúan mucho. La ubicación de la antena decide más que cualquier ajuste del programa.}

\Paso{Sitúe la antena}
{La ventana con más cielo despejado; la antena de 1090 MHz y su base magnética.}
{Vertical. Si hay una superficie metálica disponible (el marco de la ventana, una bandeja), apóyela encima.}
{Las antenas de los aviones emiten en polarización vertical; la superficie metálica actúa como plano de tierra de la antena.}
{La antena queda estable y vertical, con el cable de 3 m hasta el puesto.}
{}
{Sin ventana practicable: la recepción en interiores es posible, pero con menos alcance (etapa 7).}

\Paso{Arme la cadena}
{Cable de la antena $\rightarrow$ filtro de 1090 MHz $\rightarrow$ puerto \textbf{ANTENNA} del HackRF.}
{Si los géneros de los conectores no encajan, use un adaptador SMA del paquete del HackRF. Apriete a mano.}
{El filtro deja pasar la banda de 1090 MHz y atenúa emisoras lejanas (FM, televisión, telefonía) que pueden saturar el receptor. Es pasivo: no necesita alimentación.}
{Cadena firme, sin holguras.}
{Siga el cable con el dedo, de la antena al HackRF: no hay ningún conector suelto.}
{}

Conecte el USB a un puerto directo del PC y anote el montaje en la \hoja{U1}.

\begin{logro}
La estación está montada: antena en la mejor posición disponible, filtro en la cadena y HackRF alimentado.
\end{logro}

% ---------------------------------------------------------------------------
\section{Etapa 2 · Comprobar Ubuntu}
\label{sec:u-ubuntu}

\Puente{Cinco órdenes de terminal confirman que el sistema, el HackRF y SDRangel están listos. Abra una terminal con \tecla{Ctrl}+\tecla{Alt}+\tecla{T} y copie cada orden sin el símbolo \texttt{\$}.}

\Comando{lsb_release -ds}
{Confirmar la versión del sistema.}
{\texttt{Ubuntu 24.04.x LTS}.}
{Otra versión.}
{Avise al instructor.}

\Comando{lsusb -d 1d50:6089}
{Comprobar que el sistema ve el HackRF en el bus USB.}
{Una línea con \texttt{ID 1d50:6089}.}
{Sin salida.}
{Otro cable o puerto; \texttt{sudo dmesg | tail} muestra los mensajes al reconectar.}

\Comando{hackrf_info}
{Identificar el HackRF \emph{sin} \texttt{sudo}: así se comprueban también los permisos.}
{\texttt{Found HackRF}, \texttt{Board ID Number: 2 (HackRF One)}, número de serie y firmware.}
{\texttt{No HackRF boards found} aunque \texttt{lsusb} lo vea: faltan permisos.}
{Siga el \anexoA{sec:a-ubuntu}. No use \texttt{sudo} para trabajar.}

\Comando{snap connections sdrangel}
{Comprobar que SDRangel tiene permiso para usar dispositivos USB.}
{Una fila \texttt{raw-usb} con una conexión (no un guion).}
{\texttt{raw-usb} sin conectar.}
{\texttt{sudo snap connect sdrangel:raw-usb}.}

\Comando{snap run sdrangel}
{Abrir SDRangel como usuario normal (también puede abrirlo desde el menú de aplicaciones).}
{La ventana principal de SDRangel.}
{La ventana no aparece.}
{Lea el mensaje de la terminal y avise al instructor.}

Registre las salidas en la \hoja{U2}.

\begin{logro}
Sistema, permisos y programa verificados: lo que falle a partir de ahora será de la señal, no de la instalación.
\end{logro}

% ---------------------------------------------------------------------------
\section{Etapa 3 · Receptor y demodulador}
\label{sec:u-receptor}

\Puente{Como en el Laboratorio 01, primero se configura el dispositivo y después el canal. Esta vez el canal no produce audio, sino mensajes.}

\Paso{Configure el HackRF en 1090 MHz}
{\tecla{Ctrl}+\tecla{R} $\rightarrow$ \textbf{HackRF} $\rightarrow$ \gui{OK}; ventana \textbf{HackRF Input}.}
{Frecuencia \valor{1\,090\,000 kHz}; \gui{SR} \valor{4\,000\,000 S/s}; decimación \valor{1}; \gui{Fc} \valor{Cen}; \gui{BBF} \valor{automático}; \gui{DC} e \gui{IQ} \valor{activados}; \gui{Bias tee} y \gui{RF amp} \valor{desactivados}; \gui{LNA} \valor{24 dB}; \gui{VGA} \valor{20 dB}.}
{Los mensajes se transmiten a 2 Mchip/s: el demodulador exige al menos 2 MSa/s y trabaja por defecto a 4 MSa/s, así que la señal no se remuestrea.}
{Todos los valores se aceptan.}
{Léalos en voz alta con su compañero.}
{El HackRF no aparece en la lista: vuelva a la etapa 2.}

\Paso{Observe el silencio}
{Botón de arranque; espectro y cascada.}
{Dos minutos sin tocar nada.}
{Para reconocer un mensaje hay que saber antes cómo se ve el ruido.}
{Un piso casi plano de 4 MHz y, con tráfico, puntos o trazos muy breves en la cascada.}
{Anote el nivel del piso en la \hoja{U3}.}
{Cuadrado rojo: reconecte el USB y arranque.}

\Paso{Añada el demodulador ADS-B}
{Icono de añadir canal $\rightarrow$ \textbf{ADS-B Demodulator} $\rightarrow$ \gui{Apply} $\rightarrow$ \gui{Close}.}
{Desplazamiento \valor{0 Hz}; \gui{SR} \valor{4M}; \gui{RF BW} \valor{2,6 MHz}; umbral de correlación \valor{7 dB}; botón \gui{S} \valor{activado}; \gui{Feed} \valor{desactivado}.}
{El umbral fija cuánto debe parecerse una ráfaga al preámbulo para intentar decodificarla. \gui{S} añade los mensajes Mode S a los ADS-B. \gui{Feed} enviaría datos a terceros.}
{Una ventana con controles arriba y una tabla vacía. Estos valores son los de fábrica de SDRangel 7.27.2.}
{}
{}

\Paso{Active estadísticas, gráfico y registro}
{Botones \gui{Display Statistics} y \gui{Display Chart}; registro con \gui{Start/stop Logging ADS-B frames to .csv File}.}
{Archivo \texttt{\textasciitilde/adsb\_<grupo>.csv}.}
{Las estadísticas cuentan tramas por tipo; el CSV guarda cada trama como evidencia reproducible.}
{En pocos minutos, filas en la tabla y contadores en marcha.}
{El total de tramas crece.}
{Sin tramas tras 10 minutos: compruebe la cadena de la etapa 1. Si todo está bien y el cielo está vacío, adelante la etapa 9 y vuelva después.}

\begin{logro}
Su estación decodifica: cada fila de la tabla nació de ráfagas de 1090 MHz que superaron una comprobación de validez.
\end{logro}

% ---------------------------------------------------------------------------
\section{Etapa 4 · De la energía a la trama}
\label{sec:u-tramas}

\Puente{¿Por qué podemos fiarnos de una fila de la tabla? Porque cada mensaje recorrió una escalera: energía, preámbulo, bits y comprobación de paridad. Un pico en el espectro, por alto que sea, se queda en el primer peldaño.}

\Paso{Lea lo que el demodulador está validando}
{\gui{Display Statistics} y la columna \emph{Correlation} de la tabla.}
{Cinco minutos de recepción.}
{Las estadísticas reparten las tramas por \emph{Downlink Format} (DF) y \emph{Type Code} (TC); la correlación indica cuánto se parecía cada preámbulo al ideal.}
{Varios DF, entre ellos 17 (ADS-B) y otros como 11, 4, 5, 20 o 21 (respuestas Mode S).}
{Anote totales y distribución en la \hoja{U4}.}
{}

Después decodifique a mano una trama real (\hoja{U5}): pasará de 28 cifras hexadecimales al indicativo de un vuelo.

\Profundice{el preámbulo de 8 $\mu$s, la modulación PPM, los formatos DF y su validación, \anexoA{sec:a-trama}}

\begin{logro}
Puede explicar por qué una trama es válida y cuál es la diferencia entre ver energía y recibir un mensaje.
\end{logro}

% ---------------------------------------------------------------------------
\section{Etapa 5 · La tabla de aeronaves}
\label{sec:u-tabla}

\Puente{Cada fila agrupa los mensajes de una dirección ICAO. Pero no todo lo que muestra una fila llegó por radio: unas columnas son datos transmitidos, otras cálculos del receptor y otras consultas a bases de datos.}

\Paso{Registre tres aeronaves}
{Tabla del demodulador.}
{Las tres filas con más tramas (\emph{Frames}).}
{Más tramas implican más campos y más confianza.}
{Dirección ICAO en todas; indicativo, altitud, velocidad o posición en algunas.}
{Para cada valor de la \hoja{U6}, indique si es transmitido (nivel 5), derivado (6) o externo (7).}
{}

Es normal que falten datos: la identificación, la altitud, la velocidad y la posición viajan en mensajes distintos, con ritmos distintos (\anexoA{sec:a-trama}).

\begin{logro}
Sabe qué columnas de la tabla son evidencia de radio y cuáles no.
\end{logro}

% ---------------------------------------------------------------------------
\section{Etapa 6 · El mapa}
\label{sec:u-mapa}

\Puente{Un avión no transmite su posición completa en cada mensaje. La codifica de forma compacta (CPR), y el receptor la reconstruye con dos mensajes o con una posición de referencia: la del propio receptor. Por eso hay que configurarla.}

\begin{importante}[Privacidad]
SDRangel puede tomar la posición del receptor del sistema operativo. No publique capturas con las coordenadas exactas de la estación: recórtelas o aleje el mapa.
\end{importante}

\Paso{Fije la posición del receptor}
{\menu{Preferences $\rightarrow$ My Position}.}
{Latitud y longitud del laboratorio con dos decimales (norte y este positivos). Anote el estado de la casilla \gui{Auto-update position}.}
{Es la referencia para reconstruir posiciones y para calcular distancias.}
{Los valores se conservan al reabrir el diálogo.}
{Difieren menos de 0,05° de la ubicación real del aula.}
{}

\Paso{Abra el mapa}
{\tecla{Ctrl}+\tecla{E} $\rightarrow$ \textbf{Map} $\rightarrow$ \gui{OK}. En la tabla, una fila con Lat/Lon $\rightarrow$ \gui{Find on Feature Map}.}
{Mapa 2D (los mosaicos se descargan de Internet).}
{}
{Aviones alrededor de su estación; al seleccionar uno se resalta su fila.}
{La posición del símbolo coincide con la de la tabla.}
{Filas sin posición: espere mensajes de posición. Mapa en blanco: sin conexión para los mosaicos.}

Complete la \hoja{U7}.

\Profundice{codificación CPR, decodificación global y local, y el error de referencia que estudiará en la etapa 9, \anexoA{sec:a-cpr}}

\begin{logro}
Tiene un mapa construido con sus propias tramas y sabe qué parte de él depende de su configuración.
\end{logro}

% ---------------------------------------------------------------------------
\section{Etapa 7 · Medir la recepción}
\label{sec:u-exp}

\Puente{¿Recibe mejor con más ganancia? Responder exige un experimento honesto: el tráfico cambia de un minuto a otro y puede engañarle. Por eso se repite la condición inicial al final (diseño A--B--A).}

\begin{longtable}{@{}>{\raggedright\arraybackslash}p{0.13\linewidth}>{\raggedright\arraybackslash}p{0.47\linewidth}>{\raggedright\arraybackslash}p{0.32\linewidth}@{}}
\toprule
\textbf{Condición} & \textbf{Ajuste (5 minutos cada una)} & \textbf{Qué compara} \\
\midrule
\endhead
A & LNA 24 / VGA 20, umbral 7 dB, con filtro & Referencia \\
B & LNA 32 / VGA 30, umbral 7 dB, con filtro & Efecto de la ganancia \\
A' & Igual que A & ¿Cambió el tráfico? \\
C & Igual que A, umbral 5 dB & Efecto del umbral \\
D & Igual que A, \emph{sin} filtro (detenga el HackRF para retirarlo) & Efecto del filtro \\
\bottomrule
\end{longtable}

Antes de cada condición, reinicie las estadísticas con clic derecho sobre \gui{Display Statistics}. Escriba la predicción y los resultados en la \hoja{U8}. Si A y A' difieren más que A y B, el tráfico cambió y el efecto de la ganancia no está demostrado: dígalo así. Vuelva a la condición A al terminar.

\begin{logro}
Puede defender con datos qué mejora su estación y qué no, sin confundirlo con cambios del tráfico.
\end{logro}

% ---------------------------------------------------------------------------
\section{Etapa 8 · El analizador en 1090 MHz}
\label{sec:u-siglent}

\Puente{Un analizador de barrido recorre las frecuencias de una en una; un mensaje dura 120 $\mu$s. Casi nunca coinciden. Con una configuración especial, el analizador puede mostrar la forma de los mensajes en el tiempo, aunque nunca su contenido.}

Detenga el HackRF y lleve la antena (sin el filtro) al \textbf{RF INPUT} del analizador con el adaptador N -- SMA.

\Paso{Barrido con Max Hold}
{\gui{Frequency}, \gui{Span}, \gui{BW}, \gui{Amplitude}, \gui{Trace}.}
{Centro \valor{1090 MHz}; span \valor{10 MHz}; RBW \valor{1 MHz}; VBW \valor{3 MHz}; referencia \valor{$-40$ dBm}; preamplificador \valor{activado}; traza \valor{Max Hold} durante 2 minutos.}
{Max Hold acumula las pocas coincidencias entre barrido y ráfaga.}
{Una elevación alrededor de 1090 MHz que crece con el tiempo; en \valor{Clear Write} la traza parece solo ruido.}
{}
{}

\Paso{Span cero con disparo de vídeo}
{\gui{Span $\rightarrow$ Zero Span}; \gui{Sweep $\rightarrow$ Sweep Time}; \gui{Trigger $\rightarrow$ Video}.}
{Centro \valor{1090 MHz}; RBW \valor{1 MHz}; tiempo de barrido \valor{mínimo} (900 $\mu$s según el manual del equipo); nivel de disparo \valor{6--10 dB} por encima del piso.}
{En span cero el eje horizontal es tiempo, y el disparo de vídeo congela la pantalla cuando llega una ráfaga.}
{Bloques de unos 64 o 120 $\mu$s de ancho: mensajes cortos y largos.}
{Mida la duración con dos marcadores en modo delta.}
{Nunca dispara: baje el nivel en pasos de 2 dB. Dispara siempre: súbalo.}

Registre en la \hoja{U9} y devuelva la cadena al HackRF.

\Profundice{por qué el barrido no ve las ráfagas y qué muestra el span cero, \anexoA{sec:a-barrido}}

\begin{logro}
Ha visto la forma temporal de los mensajes de 1090 MHz y sabe por qué el analizador no puede decodificarlos.
\end{logro}

% ---------------------------------------------------------------------------
\section{Etapa 9 · Contraste con una grabación de referencia}
\label{sec:u-offline}

\Puente{¿Cómo saber si una estación decodifica \emph{bien}? Comparándola con datos de contenido conocido. La carpeta \texttt{datos\_adsb} contiene una grabación real de 1090 MHz, verificada con tres decodificadores, de un único avión volando cerca de Sicilia. Úsela también para descubrir un error clásico.}

\begin{prediccion}
Si reproduce una grabación hecha en Italia con la posición del receptor configurada en su aula, ¿dónde cree que aparecerá el avión en el mapa?
\end{prediccion}

\Paso{Cargue la grabación}
{Detenga \textbf{HackRF Input}. \tecla{Ctrl}+\tecla{R} $\rightarrow$ \textbf{FileInput}; botón de abrir archivo.}
{\texttt{\textasciitilde/datos\_adsb/modes1\_dump1090\_2MSps\_1090MHz.sdriq}; bucle \valor{activado}; aceleración \valor{1}.}
{La cabecera del archivo aporta la tasa y la frecuencia.}
{\valor{2000 kS/s}, \valor{1\,090\,000 kHz}, duración \valor{00:00:00.178}, fecha de inicio 05-01-2013.}
{}
{El archivo no abre: compruebe que está en su carpeta personal (\anexoA{sec:a-grabacion}).}

\Paso{Reprodúzcala con su posición actual}
{Icono de añadir canal de \textbf{FileInput} $\rightarrow$ \textbf{ADS-B Demodulator}; botón de reproducción.}
{Desplazamiento \valor{0 Hz}; \gui{SR} \valor{2M}; umbral \valor{7 dB}; \gui{S} activado. \gui{My Position} sin cambios.}
{2M coincide con la tasa del archivo.}
{Una fila \textbf{4d2023}, indicativo \textbf{AMC421}, y una posición en el mapa.}
{Anote esa posición en la \hoja{U10}. ¿Está cerca de Sicilia?}
{}

\Paso{Repita con una referencia correcta}
{\menu{Preferences $\rightarrow$ My Position}.}
{Anote sus valores. Desactive \gui{Auto-update position}. Introduzca latitud \valor{37,0} y longitud \valor{14,0}.}
{La posición se reconstruye eligiendo la zona más cercana a la referencia; con la referencia lejos, se elige mal.}
{La aeronave aparece junto a la costa sur de Sicilia, hacia 36,96--37,10°N y 13,79--13,86°E.}
{Compare con los valores de referencia de la \hoja{U10}.}
{Los valores vuelven solos: \gui{Auto-update position} sigue activada.}

\Paso{Restaure su posición}
{\menu{Preferences $\rightarrow$ My Position}.}
{Los valores anotados y el estado anterior de \gui{Auto-update position}.}
{Olvidarlo falsearía posiciones y distancias en vivo.}
{}
{Reabra el diálogo y compruebe.}
{}

\begin{importante}[Etiquete el origen]
Toda captura de esta etapa se marca \textbf{GRABACIÓN}. El archivo tiene el tiempo comprimido y se repite en bucle: las tasas de tramas y el alcance no describen el tráfico real (\anexoA{sec:a-grabacion}).
\end{importante}

\begin{logro}
Ha comprobado su cadena con datos conocidos y ha visto, con sus propios ojos, que una posición «decodificada» puede ser falsa sin que nada lo advierta.
\end{logro}

% ---------------------------------------------------------------------------
\section{Etapa 10 · Cierre y entrega}
\label{sec:u-informe}

\begin{enumerate}
  \item Detenga y cierre SDRangel. Desconecte el USB y después la cadena de antena.
  \item Apague el analizador; guarde el filtro y los adaptadores.
  \item Complete la interpretación final (\hoja{U11}) y reúna los entregables de la lista del Anexo B (p.~\pageref{B-sec:b-entregau}).
  \item Responda el cuestionario final (Anexo B, p.~\pageref{B-sec:b-finalu}).
\end{enumerate}

\begin{logro}[Laboratorio 02 completado]
Ha construido una estación de vigilancia pasiva y puede justificar cada dato que muestra: qué fue recibido, qué fue calculado, qué fue consultado y qué fue grabado.
\end{logro}
